% Bibitems only: insert inside the parent book's thebibliography.
\bibitem{PhysFominZelevinsky}
S. Fomin and A. Zelevinsky, \emph{Total positivity: tests and parametrizations},
Math. Intelligencer \textbf{22} (2000), 23--33;
\href{https://arxiv.org/abs/math/9912128}{arXiv:math/9912128}.

\bibitem{PhysOsterwalderSchrader}
K. Osterwalder and R. Schrader, \emph{Axioms for Euclidean Green's functions},
Commun. Math. Phys. \textbf{31} (1973), 83--112;
\href{https://doi.org/10.1007/BF01645738}{doi:10.1007/BF01645738}.

\bibitem{PhysElliottEvans}
G. A. Elliott and D. E. Evans, \emph{The structure of the irrational rotation
$C^*$-algebra}, Ann. Math. \textbf{138} (1993), 477--501;
\href{https://doi.org/10.2307/2946556}{doi:10.2307/2946556}.

\bibitem{PhysJacobson}
T. Jacobson, \emph{Entanglement equilibrium and the Einstein equation},
Phys. Rev. Lett. \textbf{116} (2016), 201101;
\href{https://arxiv.org/abs/1505.04753}{arXiv:1505.04753}.

\bibitem{PhysStinespring}
W. F. Stinespring, \emph{Positive functions on $C^*$-algebras},
Proc. Amer. Math. Soc. \textbf{6} (1955), 211--216.

\bibitem{PhysCong}
I. Cong, S. Choi and M. D. Lukin, \emph{Quantum convolutional neural networks},
Nature Physics \textbf{15} (2019), 1273--1278;
\href{https://arxiv.org/abs/1810.03787}{arXiv:1810.03787}.

\bibitem{PhysZheng}
J. Zheng, Q. Gao and Y. Lv, \emph{Quantum graph convolutional neural networks},
\href{https://arxiv.org/abs/2107.03257}{arXiv:2107.03257} (2021).

\bibitem{PhysNing}
A. Ning, T. Y. Li and N. Y. Chen, \emph{Quantum graph attention network:
A novel quantum multi-head attention mechanism for graph learning},
\href{https://arxiv.org/abs/2508.17630}{arXiv:2508.17630} (2025).

\bibitem{PhysRamsauer}
H. Ramsauer et al., \emph{Hopfield networks is all you need}, ICLR 2021;
\href{https://arxiv.org/abs/2008.02217}{arXiv:2008.02217}.

\bibitem{PhysPastawski}
F. Pastawski, B. Yoshida, D. Harlow and J. Preskill, \emph{Holographic quantum
error-correcting codes: toy models for the bulk/boundary correspondence},
JHEP \textbf{2015} (6), 149;
\href{https://arxiv.org/abs/1503.06237}{arXiv:1503.06237}.

\bibitem{PhysShadows}
H.-Y. Huang, R. Kueng and J. Preskill, \emph{Predicting many properties of a
quantum system from very few measurements}, Nature Physics \textbf{16}
(2020), 1050--1057;
\href{https://arxiv.org/abs/2002.08953}{arXiv:2002.08953}.

\bibitem{PhysRosaleny}
L. E. Rosaleny et al., \emph{Peptides as versatile scaffolds for quantum computing},
\href{https://arxiv.org/abs/1708.09440}{arXiv:1708.09440} (2017),
including supplementary information S2. Coherence observations cited here
refer to the frozen-solution EPR measurements, not a calibrated QPU.

\bibitem{PhysNitz}
M. Nitz, M. Sherawat, K. J. Franz, E. Peisach, K. N. Allen and B. Imperiali,
\emph{Structural origin of the high affinity of a chemically evolved
lanthanide-binding peptide}, Angew. Chem. Int. Ed. \textbf{43} (2004),
3682--3685; \href{https://doi.org/10.1002/anie.200460028}{doi:10.1002/anie.200460028}.
Structure: \href{https://www.rcsb.org/structure/1TJB}{PDB 1TJB}.

\bibitem{PhysRFD3}
J. Butcher et al., \emph{De novo design of all-atom biomolecular interactions
with RFdiffusion3}, bioRxiv 2025.09.18.676967 (2025);
\href{https://pmc.ncbi.nlm.nih.gov/articles/PMC12458353/}{primary preprint}.
Implementation documentation:
\href{https://github.com/RosettaCommons/foundry/blob/production/models/rfd3/README.md}{RosettaCommons Foundry RFD3}.

\bibitem{PhysThermoGFN}
A. Schreiber, \emph{ThermoGFN-IF for catalysis: a protein sequence design model
tuned with GFlowNets for stable protein design and kinetic-aware enzyme
engineering}, author project description (10 March 2026),
\href{https://huggingface.co/blog/AmelieSchreiber/thermogfn-if}{project description};
\href{https://github.com/amelie-iska/ThermoGFN-IF}{implementation scaffold}.
Cited as a proposed methodology, not completed molecular or qubit validation.

\bibitem{PhysSPURS}
\emph{Generalizable and scalable protein stability prediction with rewired
protein generative models}, Nature Communications (2026),
\href{https://pmc.ncbi.nlm.nih.gov/articles/PMC12830971/}{primary article},
\href{https://doi.org/10.1038/s41467-025-67609-4}{doi:10.1038/s41467-025-67609-4};
correction \href{https://doi.org/10.1038/s41467-026-70929-8}{doi:10.1038/s41467-026-70929-8}.

\bibitem{PhysBioEmu}
S. Lewis et al., \emph{Scalable emulation of protein equilibrium ensembles
with generative deep learning}, Science (2025), eadv9817;
\href{https://doi.org/10.1126/science.adv9817}{doi:10.1126/science.adv9817}.
Scope of the cited v1 model:
\href{https://github.com/microsoft/bioemu/blob/main/MODEL_CARD.md}{official BioEmu model card}.

\bibitem{PhysUMA}
B. M. Wood et al., \emph{UMA: A family of universal models for atoms},
\href{https://arxiv.org/abs/2506.23971}{arXiv:2506.23971} (2025).

\bibitem{PhysBreuckmannTerhal}
N. P. Breuckmann and B. M. Terhal, \emph{Constructions and noise threshold
of hyperbolic surface codes}, IEEE Trans. Inform. Theory \textbf{62} (2016),
3731--3744; \href{https://arxiv.org/abs/1506.04029v2}{arXiv:1506.04029v2}.

\bibitem{PhysDelfosse}
N. Delfosse, \emph{Tradeoffs for reliable quantum information storage in
surface codes and color codes}, IEEE ISIT (2013), 917--921;
\href{https://arxiv.org/abs/1301.6588}{arXiv:1301.6588}.
